\documentclass[11pt,a4paper]{article}
\usepackage[utf8]{inputenc}
\usepackage[T1]{fontenc}
\usepackage[french]{babel}
\usepackage{amsmath, amsthm, amssymb}
\usepackage{geometry}
\geometry{margin=1in}

\title{Le Problème des Distances Unités d'Erdős : Incidences et Déconstruction Algébrique}
\author{}
\date{}

\begin{document}
\maketitle

\begin{abstract}
Ce document expose une approche structurelle visant à résoudre le problème des distances unités d'Erdős. Nous formalisons la base axiomatique des incidences points-cercles et proposons un nouveau lemme de décomposition utilisant la topologie algébrique et la méthode polynomiale.
\end{abstract}
\vfill
\noindent Charles EDOU NZE, chercheur indépendant
\newpage

\section{Introduction et Définitions Axiomatiques}
Le problème des distances unités d'Erdős, formulé en 1946, demande quel est le nombre maximum de distances unités pouvant exister parmi $n$ points dans le plan euclidien. Erdős a prouvé une borne inférieure de $n^{1+c/\log\log n}$ et une borne supérieure de $\mathcal{O}(n^{3/2})$, qui a ensuite été améliorée par Szemerédi et Trotter.

\textbf{Définition 1 (Graphe de Distances Unités).}
Soit $P \subset \mathbb{R}^2$ un ensemble fini de points tel que $|P| = n$. Le graphe de distances unités $G = (P, E)$ est défini en posant $\{p, q\} \in E$ si et seulement si $||p - q||_2 = 1$. Le nombre total d'arêtes est noté $u(P) = |E|$.

\textbf{Définition 2 (Fonction de Distances Unités).}
Le nombre maximal de distances unités pour tout ensemble de $n$ points est donné par $u(n) = \max_{|P|=n} u(P)$.

\textbf{Conjecture (Erdős, 1946).}
Pour tout $\varepsilon > 0$, il existe une constante $C_\varepsilon$ telle que $u(n) \le C_\varepsilon n^{1+\varepsilon}$.

\section{Recherche de Littérature Contextuelle}
Le résultat fondamental de Szemerédi et Trotter (1983) borne le nombre d'incidences point-courbe. En considérant chaque point comme le centre d'un cercle unité, le problème se traduit par la limitation des incidences entre $n$ points et $n$ cercles unités.
L'application de la méthode polynomiale par Guth et Katz (2015) au problème des distances distinctes d'Erdős a introduit de nouveaux cadres algébriques.

\section{Stratégie de Preuve et Lemmes}
Notre stratégie consiste à décomposer le graphe des distances unités en sous-graphes bipartites à complexité algébrique contrôlée.

\textbf{Lemme 1 (Partitionnement Polynomial pour Cercles).}
Pour tout ensemble de $n$ points dans $\mathbb{R}^2$ et un entier $D$, il existe un polynôme non nul $F \in \mathbb{R}[x,y]$ de degré au plus $D$ tel que $\mathbb{R}^2 \setminus Z(F)$ est partitionné en $\mathcal{O}(D^2)$ cellules, chacune contenant au plus $\mathcal{O}(n/D^2)$ points.

\textit{Stratégie de preuve:} Ceci étend le partitionnement de Guth-Katz en considérant spécifiquement les intersections avec des variétés représentant des cercles unités.

\textbf{Lemme 2 (Borne d'Incidence sur les Courbes Algébriques).}
Le nombre d'incidences entre $n$ points et $m$ cercles unités, sous la condition qu'au plus $C$ cercles s'intersectent en une paire commune de points, est strictement borné par $\mathcal{O}(m^{2/3}n^{2/3} + m + n)$.

\textit{Stratégie de preuve:} Nous appliquons les bornes d'incidence dérivées des nombres de croisements et de Szemerédi-Trotter.

\section{Preuve Informelle (Zéro Ellipse)}

Détaillons explicitement l'application du Lemme 2.
Soit $P$ un ensemble de $n$ points. Pour chaque $p \in P$, définissons un cercle unité $C_p = \{x \in \mathbb{R}^2 \mid ||x - p||_2 = 1\}$.
Soit $\mathcal{C} = \{C_p \mid p \in P\}$. Ainsi, nous avons $|\mathcal{C}| = n$ cercles unités.
Le nombre de distances unités $u(P)$ est exactement la moitié du nombre d'incidences entre les points $P$ et les cercles $\mathcal{C}$. Formellement, $2 u(P) = I(P, \mathcal{C}) = |\{(p, C_q) \in P \times \mathcal{C} \mid p \in C_q\}|$.

Considérons le graphe d'incidence $H$. Pour deux cercles distincts $C_p, C_q \in \mathcal{C}$, leur intersection $|C_p \cap C_q|$ est au plus de 2.
Il n'y a donc pas de $K_{2,3}$ dans le graphe d'incidence.
Appliquons le théorème de Kővári-Sós-Turán.
Pour chaque point $p_i \in P$, soit $d_i$ son degré (le nombre de cercles sur lesquels il se trouve).
Le nombre de paires de cercles contenant un point commun est $\sum_{i=1}^n \binom{d_i}{2}$.
Puisque chaque paire de cercles s'intersecte au plus deux fois, le nombre maximal de paires de cercles partageant un point sur l'ensemble de $P$ est borné par $2 \binom{n}{2} = n(n-1)$.
Ainsi, $\sum_{i=1}^n \frac{d_i(d_i - 1)}{2} \le n(n-1)$.
En utilisant l'inégalité de Cauchy-Schwarz, $\sum_{i=1}^n d_i^2 \ge \frac{1}{n} (\sum_{i=1}^n d_i)^2$.
Par conséquent, $\frac{1}{2n} (\sum d_i)^2 - \frac{1}{2} \sum d_i \le n^2$.
Soit $I = \sum d_i$. Nous avons $\frac{I^2}{2n} - \frac{I}{2} \le n^2$.
La résolution de cette inéquation quadratique donne $I \le \sqrt{2n^3} + \frac{n}{2}$.
Ainsi, $u(P) = I/2 \le \frac{1}{2}\sqrt{2} n^{3/2} + \frac{n}{4}$.
Cela établit la borne en $O(n^{3/2})$ directement, sans aucune omission. $\blacksquare$

\section{Architecture d'Autoformalisation (Lean 4)}
\begin{verbatim}
import Mathlib.Data.Real.Basic
import Mathlib.Analysis.InnerProductSpace.PiL2
import Mathlib.Combinatorics.SimpleGraph.Basic

/-!
# Architecture formelle - Le Probleme des Distances Unites d'Erdos
-/

abbrev Point := EuclideanSpace Real (Fin 2)

def is_unit_distance (p q : Point) : Prop :=
  dist p q = 1

def UnitDistanceGraph (P : Finset Point) : SimpleGraph Point where
  Adj p q := p \in P /\ q \in P /\ p != q /\ is_unit_distance p q
  symm := by
    intro p q h
    sorry
  loopless := by
    intro p h
    sorry

def ErdosUnitDistanceConjecture : Prop :=
  forall epsilon > 0, exists C : Real, C > 0 /\
  forall P : Finset Point,
  (UnitDistanceGraph P).edgeFinset.card <= C * (P.card : Real) ^ (1 + epsilon)

lemma point_circle_incidence_bound (P : Finset Point) :
  (UnitDistanceGraph P).edgeFinset.card <= (2 ^ (1/2)) * (P.card : Real) ^ (3/2) := by
  sorry
\end{verbatim}

\end{document}
